\section{问题重述}

\subsection{问题背景}

地球同步轨道带距地面约 35786~km，通信、导航增强、气象和预警等高价值卫星大多部署于此，这一区域的空间态势感知因而受到持续关注。地基光学和雷达设备受距离和大气条件限制，难以获得高轨目标的细节信息。由观测卫星机动到目标附近、在合适的距离和光照条件下近距离成像，是获取高轨目标精细特征的主要手段，这种方式称为抵近观测。美国太空军的 GSSAP 卫星就运行在同步带附近，具备交会与近距离操作能力，用于跟踪同步带上的人造目标并识别其特性\cite{gssap}。

抵近观测对相对几何有较严格的要求。观测卫星与目标的距离既不能太远，否则成像分辨率不足，也不能太近，否则存在碰撞风险。光照条件和观测方向同样受限，观测卫星需要处在目标被太阳照亮的一侧，并以地球为背景进行观测，这样才能获得清晰的图像。上述条件还必须连续满足一段时间，才能完成一次有效观测。张海涛等针对单个地球静止轨道目标研究了绕飞持续观测的轨迹规划，使成像卫星在整个绕飞过程中保持有利的光照和观测角\cite{zhang2024}。

当观测对象从单个目标扩展到多个目标时，问题还要决定访问的顺序和时间。Alfriend 等研究了同步轨道卫星在轨服务的最优规划\cite{alfriend2006}，Zhou 等研究了多颗同步轨道卫星加注任务的规划优化\cite{zhou2015}，王峥嵘等提出了同步轨道多目标飞掠成像的任务规划方法\cite{wang2025}。这类问题把离散的排序决策和连续的轨道机动耦合在一起，求解难度随目标数量迅速增加。另一方面，高轨卫星携带的推进剂有限，每一次轨道机动都会消耗卫星的在轨寿命。因此，如何以尽量少的速度增量依次抵近多颗目标并完成观测，是高轨抵近观测任务规划的核心问题。

\subsection{问题提出}

题目给出红方观测卫星 R1 和蓝方 30 颗目标卫星 B1 至 B30 在初始历元 2027 年 1 月 1 日 00:00:00 UTC 的位置和速度，坐标系为地心惯性系 ICRF。目标卫星均位于地球同步轨道带附近，全程不机动。轨道运动采用地球中心引力二体模型，引力常数 $\mu = 398600.4418~\mathrm{km^3/s^2}$，不考虑 $J_2$ 项、日月引力和光压等摄动，太阳只用于确定光照方向。

一次有效观测需要在至少 60~s 内连续满足三个条件。第一个条件是观测卫星与目标的距离保持在 10~km 至 50~km 之间。第二个条件是光照角不超过 $45^\circ$，光照角是以目标为顶点、指向太阳的射线与指向观测卫星的射线之间的夹角。第三个条件是侦测角不超过 $45^\circ$，侦测角是以目标为顶点、指向地心的射线与指向观测卫星的射线之间的夹角。两个角度都按有向射线计算，取值范围为 $0^\circ$ 至 $180^\circ$。观测卫星采用瞬时脉冲机动，每完成一次有效观测之前最多机动两次，单次速度增量不超过 30~m/s，同一颗卫星相邻两次机动的间隔不少于 3600~s，每颗观测卫星的累计速度增量不超过 300~m/s。在上述条件下，题目提出了以下三个问题。

{\heiti 问题一}\quad 要求在初始历元后 48~h 内，通过 R1 的轨道机动完成对目标 B1 的一次有效观测。方案需要给出每次机动的时刻和三维速度增量、各段自由飞行时间以及观测的起止时刻，并验证全部观测约束，使累计速度增量最小。

{\heiti 问题二}\quad 将任务扩展到 30 天，要求 R1 依次机动，对 30 颗目标尽可能多地完成有效观测。优化目标分为两级，首先使成功观测的不同目标数量最多，在此基础上再使 R1 的累计速度增量最小。

{\heiti 问题三}\quad 为弥补单颗卫星观测能力的不足，红方提前发射了一颗偏心率小于 0.001 的高轨卫星 R2，与 R1 一同在 30 天内观测这 30 颗目标。R2 在 $t_0$ 时刻的轨道由我方设计。R2 的机动约束和速度增量预算与 R1 相同，两颗卫星的预算相互独立。优化目标仍按两级排序，首先使两颗卫星合计观测的不同目标数量最多，再使两颗卫星的累计速度增量之和最小。
